% ID: SEL-D0031-OA-E01
% ESTADO: cotejada. Segunda pasada visual completa y borrador compilado revisados.
% Fuente: Selectividad/matematicas_aplicadas/titular_junio_6 EXAMEN MAT APLIC C S.pdf
% SHA-256: bebc339cfd8f69c7c4fae160fd165eaef5f6bc93cfd6d15cb9ce289b7029d167
% Página: 1; opción A; ejercicio original 1.

Las filas de la matriz $P$ indican los respectivos precios de tres artículos $A_1$, $A_2$ y $A_3$ en dos comercios, $C_1$ (fila 1) y $C_2$ (fila 2): $P=\begin{pmatrix}25&20&15\\23&25&17\end{pmatrix}$.
Cati desea comprar 2 unidades del artículo $A_1$, 1 de $A_2$ y 3 de $A_3$.

Manuel desea comprar 5 unidades de $A_1$, 1 de $A_2$ y 1 de $A_3$.

Han dispuesto esas compras en la matriz $Q$: $Q=\begin{pmatrix}2&1&3\\5&1&1\end{pmatrix}$
\begin{enumerate}[label=\alph*)]
\item \textbf{(1.8 puntos)} Calcule $P\cdot Q^t$ y $Q\cdot P^t$ e indique el significado de los elementos de las matrices resultantes.
\item \textbf{(0.7 puntos)} A la vista de lo obtenido en el apartado anterior, \textquestiondown{}dónde les interesa hacer la compra a cada uno?
\end{enumerate}
